\chapter{Ретроспективный мост от наивного атласа к параметру порядка}

Ранний атлас зафиксировал одиннадцать коротких $\pi$-соотношений между
параметрами Стандартной модели. Последующие коллективный и
спектрально-адресный аудиты показали, что эти формулы обладают общей
арифметической организацией, однако её операторный источник оставался
неопределённым.

Том VI даёт новый мотив для возвращения к этим соотношениям: независимо
получен точный контрольный спектр проекторной фазы. Поэтому дроби атласа
проверяются как возможные собственные значения и степени щели единого
параметра порядка. Целью является восстановление их структурного
происхождения, а не отдельный подбор коэффициентов для каждой
наблюдаемой.

\section{Точный контрольный спектр Тома VI}

В минимальном энтропийно-ориентационном функционале при
$\kappa_c=\log4$ сосуществуют изотропная и одноосная фазы. Спектр
упорядоченного состояния равен
\begin{equation}
 R_{\rm crit}=\operatorname{diag}
 \left(\frac23,\frac16,\frac16\right).
 \label{eq:v6-atlas-critical-R}
\end{equation}
Для
\begin{equation}
 Q_{\rm crit}=R_{\rm crit}-\frac13I_3
\end{equation}
и щели $\Delta_{\rm crit}=1/2$ получаем
\begin{equation}
 \Spec Q_{\rm crit}
 =\left(\frac13,-\frac16,-\frac16\right),
 \qquad
 \Spec\frac{Q_{\rm crit}}{\Delta_{\rm crit}}
 =\left(\frac23,-\frac13,-\frac13\right).
 \label{eq:v6-atlas-normalized-Q}
\end{equation}

Именно нормированный оператор $Q/\Delta$ позднее вошёл в проверку Каллиаса.
Поэтому его рациональный спектр является уже существующим проектным
объектом, а не новой подстановкой ради атласа.

\section{Неожиданная таблица рангов двадцать четыре}

Старый SU(5)-аудит получил каноническое разложение
\begin{equation}
 24=8_C+3_W+1_Y+6_X+6_{\overline X}.
 \label{eq:v6-atlas-su5-ranks}
\end{equation}
Теперь возникают точные тождества
\begin{align}
 \left(\frac23,-\frac13,-\frac13\right)
 &=\left(\frac{24-8}{24},-\frac8{24},-\frac8{24}\right),
 \label{eq:v6-atlas-Q-rank-identity}\\
 \left(\frac23,\frac16,\frac16\right)
 &=\left(\frac{16}{24},\frac4{24},\frac4{24}\right),
 \qquad 16=24-8,\qquad 4=3+1,
 \label{eq:v6-atlas-R-rank-identity}\\
 \Delta_{\rm crit}&=\frac12=\frac{12}{24},\\
 \Delta_{\rm crit}^2&=\frac14=\frac6{24}.
 \label{eq:v6-atlas-gap-rank-identity}
\end{align}
Последняя строка эквивалентна целочисленному равенству
\begin{equation}
 12^2=24\cdot6.
\end{equation}

Полный перестановочный контроль даёт единственный набор ролей с точностью
до очевидных симметрий:
\begin{enumerate}
 \item ранг восемь выбирает только цветовой блок $C$;
 \item ранг четыре получается только как $W+Y=3+1$;
 \item ранг шесть выбирает $X$ или зарядово-сопряжённый $\overline X$.
\end{enumerate}

Это сильнее общего замечания ``все дроби делят 24''. Одни и те же
физические блоки восстанавливают полный спектр нормированной массы,
контрольный спектр состояния и квадрат щели.

Однако формула $(16,4,4)/24$ не является буквальным разложением
присоединённого пространства на три непересекающихся блока: ранг четыре
использован дважды, а ранг шестнадцать уже содержит один $W+Y$. Это
спектральное тождество. Повтор второй четвёрки должен быть выведен
отдельной кратностью, а не объявлен новым сектором.

\section{Словарь старых дробей}

Обозначим собственные значения нормированной массы через
\begin{equation}
 q_+=\frac23,
 \qquad q_-=-\frac13.
\end{equation}
Тогда повторяющиеся коэффициенты наивного атласа получают единый словарь:
\begin{align}
 \frac13&=-q_-,&
 \frac23&=q_+,&
 \frac43&=1-q_-,\\
 \frac32&=q_+^{-1},&
 \frac12&=\Delta_{\rm crit},&
 \frac14&=\Delta_{\rm crit}^2.
 \label{eq:v6-atlas-fraction-dictionary}
\end{align}

Семь из одиннадцати коротких строк прежнего $\pi$-атласа переписываются
без остатка:
\begin{align}
 \alpha_s
 &=\frac1{6+\pi^2\Delta_{\rm crit}^2},\label{eq:v6-atlas-alpha-rewrite}\\
 \sin^2\theta_W
 &=\frac{8-3\Delta_{\rm crit}^2/\pi}
 {(24-3)+(3+1)\pi},\label{eq:v6-atlas-weinberg-rewrite}\\
 \frac{m_b}{m_p}&=\pi+1-q_-,\\
 \left(\frac{m_s}{m_p}\right)^{-1}&=\pi^2-q_-,\\
 \left(\frac{m_\tau}{m_\mu}\right)_{\rm core}
 &=(\pi+1)^2+q_-,\\
 \Omega_{dm}&=\pi^{-1}-\Delta_{\rm crit}\pi^{-2},\\
 \Omega_b&=\Delta_{\rm crit}\pi^{-2}.
 \label{eq:v6-atlas-seven-rewrites}
\end{align}
Максимальный машинный остаток этих семи тождеств равен нулю в double-
арифметике. Дополнительно ведущая дробь мюонной формулы равна
$3/2=q_+^{-1}$, а поправка тау $2\alpha/3=\alpha q_+$.

\begin{proposition}[От свободных дробей к спектру порядка]
Значительная часть малых рациональных коэффициентов раннего атласа может
быть прочитана как спектр $Q_{\rm crit}/\Delta_{\rm crit}$ и степени его
щели. Поэтому прежний оператор спектральных адресов естественно расширить
от одной дилатации $\Pi$ до совместного исчисления функций
$\Pi$, $R_{\rm crit}$ и $Q_{\rm crit}/\Delta_{\rm crit}$.
\end{proposition}

\section{Главный контроль против самообмана}

Точная таблица относится к контрольному переходу при $\kappa=\log4$.
Поздний канонический термический родитель при
$\beta_c=1.5426695\ldots$ имеет другой спектр:
\begin{equation}
 (0.9121666\ldots,0.0439167\ldots,0.0439167\ldots)
\end{equation}
и щель $0.8682499\ldots$. Максимальное отличие от
\eqref{eq:v6-atlas-critical-R} равно $0.2454999\ldots$, а квадрат новой
щели равен $0.7538579\ldots$, не $1/4$.

Следовательно, ранговый 24 мост не описывает автоматически текущий
динамический минимум. Возможны только два честных чтения:
\begin{enumerate}
 \item спектр $(2/3,1/6,1/6)$ является особой критической контрольной
 точкой или промежуточным универсальным состоянием;
 \item совпадение является ретроспективной арифметической подсказкой,
 которая исчезает после выбора полного родительского функционала.
\end{enumerate}

\begin{nogocriterion}[Запрет возврата к статусу атласа]
Точные переписывания не выводят массы, связи или космологические доли.
Они не выбирают наблюдаемую, масштаб перенормировки и функцию от
спектральных операторов. Пока общий functional calculus не выдаёт новую
скрытую строку до сравнения с данными, результат остаётся
ретроспективным.
\end{nogocriterion}

\section{Подсказка для текущей проверки носителя}

Наиболее содержательное равенство имеет вид
\begin{equation}
 \boxed{
 \Delta_{\rm crit}^2=\frac14=\frac{\rank X}{24}.}
 \label{eq:v6-atlas-two-copy-clue}
\end{equation}
Левая часть естественно появляется в тензорном квадрате щели параметра
порядка, правая --- как нормированный ранг смешанного цвето-слабого блока.
Одновременно спектр $(16,4,4)/24$ требует повторения ранга четыре
электрослабого адреса.

Это непосредственно уточняет следующую проверку спинорного накрытия: двухкопийный
носитель должен не просто иметь размерность два, а функториально
реализовать квадрат щели и повтор адреса ранга четыре, сохраняя калибровочный
коммутант. Если тензорный квадрат даёт лишь виртуальное чтение или смешанный
сектор с сопряжёнными зарядами, подсказка не становится физическим
носителем.

Совпадение числа $16$ с размерностью поколения $H_{16}$ также записано,
но не повышается в статусе: число шестнадцать уже имеет несколько
независимых источников, а переход $H_{15}\to H_{16}$ ранее менял родителя,
след и физические рёбра.

\begin{protocol}
\begin{enumerate}
 \item исходный атлас: \textbf{прочитан как архив, не переписан};
 \item прежние коллективный и SU(5)-аудиты: \textbf{учтены};
 \item точный спектр контрольной фазы: \textbf{$(2/3,1/6,1/6)$};
 \item нормированная масса: \textbf{$(2/3,-1/3,-1/3)$};
 \item единственный ранговый 24 словарь: \textbf{$8_C$, $3_W+1_Y$, $6_X$};
 \item равенство щели: \textbf{$\Delta^2=6/24=1/4$};
 \item точно переписанных коротких строк атласа: \textbf{семь из одиннадцати};
 \item selector наблюдаемых: \textbf{не выведен};
 \item совпадение с текущим динамическим минимумом: \textbf{нет};
 \item новая подсказка: \textbf{щель тензорного квадрата и повтор ранга четыре адреса};
 \item следующая проверка: \textbf{двухкопийная кратность спинорного накрытия}.
\end{enumerate}
\end{protocol}

Машинный сертификат сохранён в
\path{s2t_v6_naive_atlas_order_parameter_rank_bridge_gate_results.json}.